%%APPEND title="Why converse excision is an additional axiom"
\textbf{Complement (every integral homology equivalence).}\ours
Under the hypotheses of the proposition, $W$ contains every functor that induces isomorphisms on integral homology, and not only the functors $A\to e$ with $A$ acyclic.

\emph{Proof.}
For $g:A\to B$ let $K'(g)$ be the category made of $A$, of $B$, of the arrows $b\to a$ given by the arrows $b\to g(a)$ of $B$, and of one further object $t$ with one arrow $t\to a$ for every object $a$ of $A$. The full subcategories $C_1$ on $A\cup B$ and $C_2=A\cup\{t\}$ are cosieves with union $K'(g)$ and intersection $A$; $C_2$ has the initial object $t$, and the retraction $C_1\to B$ which is $g$ on $A$ has the inclusion of $B$ as left adjoint. If $K'(g)\to e$ is in $W$, then $C_2\to K'(g)$ is in $W$, (A6) gives $A\to C_1$ in $W$, and $g\in W$. The classifying space of $K'(g)$ is the mapping cone of $Bg$. If $g$ is an integral homology equivalence, $C=K'(g)$ is connected and acyclic, and $K'(C\to e)$ realizes the suspension of $BC$, which is weakly contractible. Applying the first step to $C\to e$ and then to $g$ gives $g\in W$.\hfill$\square$

So the functors inducing isomorphisms on integral homology form the least class satisfying the hypotheses of the proposition (they satisfy them by the exact sequence (A7)); the commentary on Chapter~XVI, ms.~27, treats the corresponding cube axiom and its incompatibility with the fibration axioms.
